\documentclass[10pt]{article}
\usepackage[T1]{fontenc}
\usepackage[utf8]{inputenc}
\usepackage[french]{babel}
\usepackage{amsmath,amssymb,mathrsfs,mathpazo}
\usepackage{geometry}
\geometry{a4paper,margin=2cm}
\usepackage{graphicx,array,tabularx,tabularray,colortbl,multirow,multicol,booktabs,enumitem,pifont,ulem,fancybox,diagbox,cancel}
\usepackage{pstricks,pstricks-add,pst-plot,pst-tree,pst-node,pst-text,pst-func,pst-math,pst-bspline,pst-3dplot,pst-eucl,pst-solides3d}
\let\pscalculate\undefined
\usepackage{pst-bezier}
\usepackage{tikz,pgfplots,tkz-tab}
\usetikzlibrary{arrows,arrows.meta,patterns,calc,positioning,trees,decorations.markings,shapes.geometric,shapes.misc,shapes.arrows,intersections,decorations.pathmorphing,decorations.pathreplacing,backgrounds,shadows,babel}
\pgfplotsset{compat=1.18}
\usepgfplotslibrary{fillbetween}
\usepackage[np]{numprint}
\usepackage{listings,makecell,pas-tableur,twoopt}
\let\barre\undefined
\usepackage[pstricks]{tablvar}
\definecolor{darkgreen}{RGB}{0,100,0}
\providecommand{\red}{\color{red}}
\providecommand{\blue}{\color{blue}}
\providecommand{\green}{\color{green}}
\providecommand{\black}{\color{black}}
\providecommand{\white}{\color{white}}
\providecommand{\gray}{\color{gray}}
\providecommand{\cyan}{\color{cyan}}
\providecommand{\magenta}{\color{magenta}}
\providecommand{\yellow}{\color{yellow}}
\providecommand{\violet}{\color{violet}}
\providecommand{\orange}{\color{orange}}
\providecommand{\euro}{\texteuro}
\providecommand{\EUR}{\texteuro}
\providecommand{\ieme}{\textsuperscript{e}}
\providecommand{\degres}{\ensuremath{^\circ}}
\providecommand{\degree}{\ensuremath{^\circ}}
\providecommand{\vv}[1]{\overrightarrow{#1}}
\providecommand{\Coord}[1]{#1}
\providecommand{\decosix}{}
\providecommand{\starredbullet}{\ensuremath{\bullet}}
\pagestyle{empty}
\setlength{\parindent}{0pt}
\usepackage{adjustbox}\DeclareUnicodeCharacter{2236}{:}\DeclareUnicodeCharacter{2190}{\ensuremath{\leftarrow}}
\providecommand{\R}{\mathbb{R}}
\providecommand{\N}{\mathbb{N}}
\providecommand{\Z}{\mathbb{Z}}
\providecommand{\Q}{\mathbb{Q}}
\providecommand{\D}{\mathbb{D}}
\providecommand{\C}{\mathbb{C}}
\providecommand{\e}{{\mathrm{e}}}
\providecommand{\vect}[1]{\overrightarrow{#1}}
\providecommand{\vectt}[1]{\overrightarrow{\mathrm{#1}}}
\providecommand{\vvt}[1]{\overrightarrow{\mathrm{#1}}}
\providecommand{\np}[2][]{{#2~#1}}
\providecommand{\vv}[1]{{\overrightarrow{#1}}}
\renewcommand{\d}{{\mathrm{d}}}
\providecommand{\degres}{{^\circ}}
\providecommand{\degre}{{^\circ}}
\providecommand{\blue}{\color{blue}}
\providecommand{\red}{\color{red}}
\providecommand{\enskip}{\quad}
\providecommand{\llbracket}{\lbrack\!\lbrack}
\providecommand{\rrbracket}{\rbrack\!\rbrack}
\providecommand{\euro}{{\text{€}}}
\providecommand{\no}{{\text{n°}}}
\providecommand{\interval}[2]{{[#1;#2]}}
\providecommand{\intervFO}[2]{{[#1;#2[}}
\providecommand{\dotfill}{\ldots}

\providecommand{\ds}{\displaystyle}
\providecommand{\textding}[1]{#1}
\providecommand{\Oij}{\(\mathrm{O};\vec{\imath},\vec{\jmath}\)}
\providecommand{\Oijk}{\(\mathrm{O};\vec{\imath},\vec{\jmath},\vec{k}\)}
\providecommand{\Ouv}{\(\mathrm{O};\vec{u},\vec{v}\)}
\providecommand{\intervFF}[2]{\ensuremath{[#1;#2]}}
\renewcommand{\euro}{\ifmmode\text{\texteuro}\else\texteuro\fi}
\DeclareUnicodeCharacter{00B0}{\ifmmode^{\circ}\else\textdegree\fi}
\DeclareUnicodeCharacter{20AC}{\ifmmode\text{\texteuro}\else\texteuro\fi}
% Énoncés : Éducation nationale. Transcription et corrigés : APMEP, Denis Vergès et contributeurs. Découpage et classement : CoopMaths, Sébastien Lozano, Éric Elter et contributeurs. Adaptation de la présentation : Maths Explorer.
% CC BY-SA : https://coopmaths.fr/www/about/
% https://coopmaths.fr/alea/static/e3c/2020/tex/e3c_2020_00_sujet16_2_cor.tex
\begin{document}
\noindent\textit{APMEP / CoopMaths. CC BY-SA. Présentation adaptée par Maths Explorer.}\par\medskip
\subsection*{Partie A}
	
	On considère la fonction polynôme du second degré \(P\) définie sur \(\mathbb{R}\) par :
	
	\[
	P(x) = x^2 - 7x + 6.
	\]
	
	1. La somme des racines est égale à 7 et leur produit à 6 : il est à peu près évident que ces racines sont 1 et 6.
	
	Sinon, on peut calculer :
	\[
	\Delta = 7^2 - 4 \times 6 = 49 - 24 = 25 = 5^2 > 0.
	\]
	Les racines sont donc :
	\[
	x_1 = \dfrac{7 + 5}{2} = 6 \quad \text{et} \quad x_2 = \dfrac{7 - 5}{2} = 1.
	\]
	
	2. On sait que le trinôme est positif (signe du coefficient 1) sauf entre les racines 1 et 6.
		\begin{center}
		\begin{tikzpicture}[double distance=2pt]
			\tkzTabInit{$x$/1,$P(x)$/1}{$-\infty$,$1$,$6$, $+\infty$}
			\tkzTabLine{,+,z,-,z, +}
			
		\end{tikzpicture}
	\end{center} 
	
	\subsection*{Partie B}
	
	\(f(x) = 2x^3 - 21x^2 + 36x\).
	
	1. La fonction polynôme est dérivable sur \(\mathbb{R}\) et sur cet ensemble :
	
	\[
	f'(x) = 3 \times 2x^2 - 2 \times 21x + 36 = 6x^2 - 42x + 36 = 6(x^2 - 7x + 6) = 6P(x).
	\]
	
	2. Comme 6 est positif, le signe de \(f'(x)\) est celui de \(P(x)\) et on a vu que ce polynôme est positif sauf sur l’intervalle \(\left]1 ; 6\right[\).\\
	\(f(1) = 2 - 21 + 36 = 17\)\\
	 \(f(6) = 2 \times 216 - 21 \times 36 + 36 \times 6 = 432 - 756 + 216 = 648 - 756 = -108\).
		\begin{center}
		\begin{tikzpicture}[double distance=2pt]
			\tkzTabInit{$x$/1,$f'(x)$/1,$f(x)$/2}{$-\infty$,$1$,$6$, $+\infty$}
			\tkzTabLine{,+,z,-,z, +}
			\tkzTabVar{-/,+/$17$,-/$-108$,+/}
		\end{tikzpicture}
	\end{center} 
	3. On sait qu’une équation de la tangente \(T\) est : \(y = f(3) + f'(3)(x - 3)\).
	
	Avec \(f(3) = 2 \times 3^3 - 21 \times 3^2 + 36 \times 3 = 54 - 189 + 108 = -27\) et
	
	\[
	f'(3) = 6 \times 3^2 - 42 \times 3 + 36 = 54 - 126 + 36 = -36,
	\]
	
	l’équation devient :
	
	\[
	y - (-27) = -36(x - 3) \quad \text{ou} \quad y = -36x + 108 - 27 \quad \text{et finalement} \quad y = -36x + 81.
	\]

\end{document}
